\documentclass[11pt,a4paper]{article}
\usepackage[utf8]{inputenc}
\usepackage[T1]{fontenc}
\usepackage[provide=*,french]{babel}
\usepackage{amsmath,amssymb,mathtools}
\usepackage[most]{tcolorbox}
\usepackage{geometry,hyperref}
\usepackage[expansion=false]{microtype}
\geometry{margin=2.5cm}
\hypersetup{colorlinks=true,linkcolor=blue!55!black}
\newtcolorbox{theorembox}{breakable,colback=blue!3,colframe=blue!60!black,title=Théorème}
\newtcolorbox{lemmabox}[1][]{breakable,colback=green!3,colframe=green!50!black,title=Lemme,#1}
\newtcolorbox{proofbox}[1][]{breakable,colback=gray!2,colframe=gray!65!black,title=Démonstration,#1}
\newtcolorbox{remarkbox}{breakable,colback=violet!3,colframe=violet!55!black,title=Remarque}
\newtcolorbox{resultbox}{breakable,colback=red!2,colframe=red!60!black,title=Conclusion}
\title{Jean Dieudonné -- Calcul infinitésimal\\
\large Résolution rigoureuse de l'exercice 23, pages 118--119}
\author{}
\date{}
\begin{document}
\maketitle

\section{Énoncé exactement transcrit}

\begin{theorembox}
Soit $h$ un entier supérieur ou égal à $1$. Montrer que, lorsque
$n\to+\infty$, il existe deux constantes $a>0$ et $b>0$ telles que
\[
a n^{\frac{4h-1}{2h}}
\leq
S_n:=\sum_{k=1}^{n-1}
\frac{\sin^2\!\left(\dfrac{\pi k^{2h}}{n}\right)}
     {\sin^2\!\left(\dfrac{\pi k}{n}\right)}
\leq
b n^{\frac{4h-1}{2h}}.
\]
Les constantes peuvent dépendre de $h$, mais non de $n$.
\end{theorembox}

Posons
\[
K_n:=\left\lfloor\left(\frac n2\right)^{1/(2h)}\right\rfloor,
\qquad \gamma:=\frac{4h-1}{2h}=2-\frac1{2h}.
\]

\section{Préliminaires}

\begin{lemmabox}[title=Symétrie]
Le terme de rang $n-k$ est égal au terme de rang $k$.
\end{lemmabox}
\begin{proofbox}
On a
\[
\sin^2\left(\frac{\pi(n-k)}n\right)
=\sin^2\left(\frac{\pi k}n\right).
\]
De plus $(n-k)^{2h}\equiv k^{2h}\pmod n$. Il existe donc $q\in\mathbb Z$
tel que $(n-k)^{2h}=k^{2h}+qn$, d'où
\[
\sin^2\left(\frac{\pi(n-k)^{2h}}n\right)
=\sin^2\left(\frac{\pi k^{2h}}n+q\pi\right)
=\sin^2\left(\frac{\pi k^{2h}}n\right).
\]
\end{proofbox}

\begin{lemmabox}[title=Deux inégalités trigonométriques]
Pour $0\leq x\leq\pi/2$, on a
\[
\frac{2x}{\pi}\leq\sin x\leq x.
\]
Si $p\geq1$ et $0<\theta\leq\pi/(2p)$, alors
\[
\frac{|\sin(p\theta)|}{\sin\theta}\leq p.
\]
\end{lemmabox}
\begin{proofbox}
La première double inégalité résulte de la concavité de $\sin$ sur
$[0,\pi/2]$ et de $\sin x\leq x$. Pour la seconde, la fonction
$x\mapsto\sin x/x$ est décroissante sur $(0,\pi/2]$; comme
$0<\theta\leq p\theta\leq\pi/2$, on obtient
$\sin(p\theta)/(p\theta)\leq\sin\theta/\theta$.
\end{proofbox}

\section{Majoration}

La symétrie donne, à un éventuel terme central près,
\[
S_n\leq2\sum_{1\leq k\leq n/2}
\frac{\sin^2(\pi k^{2h}/n)}{\sin^2(\pi k/n)}.
\]
Nous découpons cette somme à $K_n$.

\subsection{Petits indices : $1\leq k\leq K_n$}

Prenons $\theta=\pi k/n$ et $p=k^{2h-1}$. Alors
$p\theta=\pi k^{2h}/n\leq\pi/2$. Le lemme donne
\[
\frac{|\sin(\pi k^{2h}/n)|}{\sin(\pi k/n)}\leq k^{2h-1}.
\]
Par conséquent,
\[
\sum_{k=1}^{K_n}
\frac{\sin^2(\pi k^{2h}/n)}{\sin^2(\pi k/n)}
\leq\sum_{k=1}^{K_n}k^{4h-2}
\leq C_hK_n^{4h-1}.
\]
Comme $K_n\asymp n^{1/(2h)}$,
\[
K_n^{4h-1}\asymp n^{(4h-1)/(2h)}=n^\gamma.
\]

\subsection{Grands indices : $K_n<k\leq n/2$}

Sur cet intervalle, $\sin(\pi k/n)\geq2k/n$, tandis que le numérateur est
au plus $1$. Ainsi
\[
\sum_{K_n<k\leq n/2}
\frac{\sin^2(\pi k^{2h}/n)}{\sin^2(\pi k/n)}
\leq\frac{n^2}{4}\sum_{k>K_n}\frac1{k^2}
\leq\frac{n^2}{4K_n}.
\]
Or $n^2/K_n\asymp n^{2-1/(2h)}=n^\gamma$. Les deux zones réunies
fournissent une constante $b_h>0$ telle que
\[
S_n\leq b_hn^\gamma
\]
pour tout $n$ assez grand.

\section{Minoration}

Nous isolons un bloc sur lequel le numérateur est uniformément séparé de zéro.
Posons
\[
L_n:=\left\lceil\left(\frac n2\right)^{1/(2h)}\right\rceil,
\qquad
M_n:=\left\lfloor\left(\frac{3n}{4}\right)^{1/(2h)}\right\rfloor.
\]
Pour $L_n\leq k\leq M_n$, on a, à des erreurs tendant vers zéro dues aux
parties entières,
\[
\frac12\leq\frac{k^{2h}}n\leq\frac34.
\]
Il existe donc $n_0(h)$ tel que, pour $n\geq n_0(h)$,
\[
\sin^2\left(\frac{\pi k^{2h}}n\right)\geq c_0,
\]
où l'on peut choisir une constante absolue $c_0>0$. D'autre part,
$\sin(\pi k/n)\leq\pi k/n$. Par conséquent,
\[
S_n\geq\frac{c_0n^2}{\pi^2}
\sum_{k=L_n}^{M_n}\frac1{k^2}.
\]
Le quotient $M_n/L_n$ tend vers $(3/2)^{1/(2h)}>1$. Il existe donc
$\delta_h>0$ tel que $M_n-L_n\geq\delta_hL_n$ pour $n$ assez grand. Comme
$k\leq M_n\leq C_hL_n$ sur ce bloc,
\[
\sum_{k=L_n}^{M_n}\frac1{k^2}
\geq\frac{M_n-L_n+1}{M_n^2}
\geq\frac{c_h}{L_n}.
\]
Ainsi
\[
S_n\geq c'_h\frac{n^2}{L_n}
\geq a_hn^{2-1/(2h)}=a_hn^\gamma.
\]

\section{Conclusion}

\begin{resultbox}
Il existe deux constantes $a_h,b_h>0$ et un entier $n_0(h)$ tels que, pour
tout $n\geq n_0(h)$,
\[
\boxed{
a_h n^{\frac{4h-1}{2h}}
\leq
\sum_{k=1}^{n-1}
\frac{\sin^2(\pi k^{2h}/n)}{\sin^2(\pi k/n)}
\leq
b_h n^{\frac{4h-1}{2h}}.}
\]
Autrement dit,
\[
S_n\asymp_h n^{(4h-1)/(2h)}.
\]
\end{resultbox}

\begin{remarkbox}
La coupure $K_n\asymp n^{1/(2h)}$ est naturelle : elle correspond à
$k^{2h}/n\asymp1$. Avant cette échelle, l'inégalité
$|\sin(p\theta)|\leq p\sin\theta$ contrôle le quotient; après cette échelle,
la majoration du numérateur par $1$ et la sommation de $n^2/k^2$ prennent le
relais. La même échelle fournit le bloc de minoration.
\end{remarkbox}

\end{document}
